Pour chacune des affirmations suivantes, indiquer si elle est vraie ou fausse.

\emph{Justifier chaque réponse. Une réponse non justifiée ne rapporte aucun point.}

\medskip

\begin{enumerate}
\item On considère la fonction $f$ définie pour tout $x \in \mathbb{R}$ par $f(x) = x \e^{(x^2)}$.

\textbf{Affirmation 1} : $f$ est concave sur $]-\infty~;~0]$.

\item Un refuge spécialisé accueille 24 rapaces numérotés de 1 à 24. Le refuge travaille avec trois vétérinaires. Afin d'administrer un soin à tous les oiseaux, le directeur veut donc former trois groupes, numérotés de 1 à 3, constitués chacun de 8 rapaces.

\textbf{Affirmation 2} : Le nombre de façons de répartir les rapaces dans les groupes 1 à 3 est :

\[\frac{24!}{(8!)^3}.\]

\item Soit $f$ une fonction définie et dérivable sur $\mathbb{R}$, solution de l'équation différentielle :

\[y' = -y + 1 \text{ sur } \mathbb{R}.\]

\textbf{Affirmation 3} : La fonction $g$ définie pour tout $x \in \mathbb{R}$ par $g(x) = f(x) - 1$ vérifie

\[g' = -g.\]

\item La suite $(u_n)$ est définie pour tout entier naturel $n$ non nul par :

\[u_n = \frac{\sin(n)}{n}.\]

\textbf{Affirmation 4} : La suite $(u_n)$ n'admet pas de limite.

\item On a tracé ci-dessous la parabole d'équation $y = -x^2 + 4x + 1$ et la droite d'équation $y = x + 1$.

\begin{center}
%Début Tikz
%\begin{tikzpicture}[scale=0.8]
%\begin{axis}[
%xmin=-1, xmax=5.5, ymin=-1.5, ymax=5.5,
%axis lines=middle,
%grid=both,
%xlabel={}, ylabel={},
%width=10cm, height=8cm,
%xtick={0,1,2,3,4,5}, ytick={-1,0,1,2,3,4,5}
%]
%\addplot[domain=-1:5.5, samples=100, thick] {-x^2+4*x+1};
%\addplot[domain=-1:5.5, samples=2, thick] {x+1};
%\addplot[domain=0:3, samples=100, fill=gray, fill opacity=0.4, draw=none] {-x^2+4*x+1} \closedcycle;
%\addplot[domain=0:3, samples=2, fill=white, fill opacity=1, draw=none] {x+1} \closedcycle;
%\end{axis}
%\end{tikzpicture}
%Fin Tikz
%Début PSTricks
\psset{unit=1cm,arrowsize=2pt 3}
\begin{pspicture*}(-1.1,-1.5)(5.5,5.5)
\psaxes[linewidth=1.25pt,labelFontSize=\scriptstyle]{->}(0,0)(-1,-1.5)(5.5,5.5)
\psplot[plotpoints=2000,linewidth=1.25pt,linecolor=red]{-0.5}{4.5}{4 x mul 1 add x dup mul sub}
\psplot[plotpoints=500,linewidth=1.25pt,linecolor=blue]{-1}{5}{x  1 add}
\pscustom[fillstyle=solid,fillcolor=lightgray]{\psplot[plotpoints=2000,linewidth=1.25pt,linecolor=red]{0}{3}{4 x mul 1 add x dup mul sub}\psline(3,4)(0,1)}
\psgrid[gridlabels=0pt,subgriddiv=2,gridwidth=0.2pt,subgridwidth=0.1pt]
\end{pspicture*}
%Fin PSTricks
\end{center}

\textbf{Affirmation 5} : L'aire du domaine grisé est égale à $6{,}5$ unités d'aire.
\end{enumerate}
